\begin{definition}
Let $N$ be normal data of height $p$, with cyclic edge identities
$f_0=e_0,f_1,\ldots,f_{p-1}$, and let $1\le j\le p-1$.  A
\emph{trunk rotation from $y_0$ to $y_j$} consists of a coloring
$\varphi_j$, a new distinguished identity $f_j$, and normal data $N'$ such
that:
\begin{itemize}
\item $f_j$ is uncolored, the old color of $f_{i+1}$ is shifted to $f_i$
      for $0\le i<j$, and every identity outside
      $\{f_0,\ldots,f_j\}$ keeps its color;
\item the first $p$ vertices of the new tree are
      $y_j,y_{j+1},\ldots,y_{p-1},y_0,\ldots,y_{j-1}$, with the
      corresponding cyclic edge identities between consecutive vertices;
\item all vertices and edge identities after the trunk retain their old
      positions, and the new closing identity is $f_{j-1\pmod p}$.
\end{itemize}
The old and new trees have the same length and vertex set, and the new
normal height remains $p$.  All indexing is on individual cycle-edge
identities, so a rotation never identifies parallel edges.
\end{definition}
